\section{Geometric realizations of Demazure--Lusztig type operators}
\label{sec:geometric_hecke}

Let $\Psi=\Psi_{\rfunc}$ be a root ideal, and set $X=\SX_{\rfunc}$.
Let $S'$ be the set of simple reflections in the Weyl group $W$ preserving
the set $\Psi\cup(-\Psi)$, and let $\Hdeg'\subset\Hdeg$ denote the
subalgebra generated by $S'$ and the polynomial subalgebra $\Z[x_1,\dots,x_n,\epsilon]$.
By Lemma~\ref{lem:symmetry_of_root_ideal}, each $s_i\in S'$ satisfies
$\rfunc(k)\neq i$ for $k\geq i+2$. We use the following two constructions
according to whether $\epsilon_i-\epsilon_{i+1}$ belongs to $\Psi$.
Let $K=T\times S$ or $G\times S$.

\subsection{The two geometric constructions}

Fix $s_i\in S'$. We first consider the case where $\rfunc(i)=i-1$ and $\rfunc(i+1)=i$, so that the
simple root $\epsilon_i-\epsilon_{i+1}$ belongs to $\Psi$.
Define a reverse Hessenberg function $\rfunc_i^-$ by setting $\rfunc_i^-(i+1)=i-1$
and $\rfunc_i^-(k)=\rfunc(k)$ for all $k\neq i+1$.
Then $\rfunc_i^-$ is weakly increasing and determines morphisms
\[
  j_i:Y_i=\SX_{\rfunc_i^-}\hookrightarrow X,
  \qquad
  \pi_i:Y_i\longrightarrow B_i=\SX_{\rfunc_i^-,i}^P.
\]
By Proposition~\ref{prop:inclusion_gysin} (applied with $m=i+1$ and $a=i$), the embedding $j_i$
realizes $Y_i$ as a smooth Cartier divisor in $X$ whose normal bundle has first Chern class
$c_1^K(N_{Y_i/X})=j_i^*(x_{i+1}-x_i-\epsilon)$.
Moreover, since $\rfunc_i^-(i)=\rfunc_i^-(i+1)=i-1$ and $\rfunc_i^-(k)=\rfunc(k)\neq i$ for $k\geq i+2$,
the function $\rfunc_i^-$ satisfies \eqref{eq:pushpull_hypotheses} with $m=i$.
Thus Proposition~\ref{prop:pushpull_projective_bundle} identifies $\pi_i$ with the equivariant
projective bundle $\PP(L_{i+1}/L_{i-1})$ over the smooth base $B_i$.
Applying $j_i^*$, $\pi_{i*}$, $\pi_i^*$, and $j_{i*}$ in that order and adding the identity,
we define the degree-zero endomorphism
\begin{equation}\label{eq:hecke_case_A}
  \mathcal S_i=\id+j_{i*}\pi_i^*\pi_{i*}j_i^*
  \qquad\text{on }H^*_K(X;\Z).
\end{equation}
For the maximal root ideal $\Psi=\Delta^+$, this divisor and forgetful correspondence coincide
with the geometric setup of Cautis and Kamnitzer \cite[Lemma~9.3]{cautis2018quantum} in equivariant $K$-theory.

Next, for our fixed $s_i\in S'$, suppose that $\rfunc(i)=\rfunc(i+1)$, which by Lemma~\ref{lem:symmetry_of_root_ideal} is equivalent to $\epsilon_i-\epsilon_{i+1}\notin\Psi$.
In this case, the condition $\rfunc(k)\neq i$ for $k\geq i+2$ implies that the function $\rfunc$ itself
satisfies \eqref{eq:pushpull_hypotheses} with $m=i$.
Hence Proposition~\ref{prop:pushpull_projective_bundle} yields a projective line bundle on the entire variety:
\[
  \pi_i:X\longrightarrow B_i=\SX_{\rfunc,i}^P.
\]
Let $\mathcal E_i=\Omega^1_{X/B_i}\otimes\C_{-1}$ denote the relative cotangent line bundle
twisted by the one-dimensional representation where $G$ acts trivially and $S$ acts with weight $-1$.
Since the relative tangent bundle is $T_{X/B_i}\cong\Hom(L_i/L_{i-1},L_{i+1}/L_i)$, its dual twisted by $\C_{-1}$ has
equivariant Euler class
\[
  e_K(\mathcal E_i)=c_1^K(\mathcal E_i)
  =c_1^K(L_i/L_{i-1})-c_1^K(L_{i+1}/L_i)-\epsilon
  =x_{i+1}-x_i-\epsilon.
\]
Multiplying the push--pull operator by this degree-two class, we define the degree-zero endomorphism
\begin{equation}\label{eq:hecke_case_B}
  \mathcal S_i=\id+e_K(\mathcal E_i)\cup\pi_i^*\pi_{i*}
  \qquad\text{on }H^*_K(X;\Z).
\end{equation}

\subsection{Identification with the polynomial operators}

Let $\Phi_K:R_K\twoheadrightarrow H_K^*(X;\Z)$ denote the presentation homomorphism.

\begin{theorem}\label{thm:geometric_hecke}
For every $s_i\in S'$ and every polynomial $f\in R_K$, the geometric operators
\eqref{eq:hecke_case_A} and \eqref{eq:hecke_case_B} satisfy
\[
  \mathcal S_i\Phi_K(f)=\Phi_K(\BS_i f).
\]
In particular, the relation ideal $I_{\rfunc,K}$ is stable under $\BS_i$.
\end{theorem}

\begin{proof}
When $\epsilon_i-\epsilon_{i+1}\in\Psi$, Theorem~\ref{thm:pushpull_divided_difference},
applied to $Y_i$, identifies $\pi_i^*\pi_{i*}$ with $\partial_i$ on polynomial representatives.
Proposition~\ref{prop:inclusion_gysin}
then identifies $j_{i*}$ with multiplication by $x_{i+1}-x_i-\epsilon$.
When $\epsilon_i-\epsilon_{i+1}\notin\Psi$, the same push--pull theorem
applies on $X$, and the Euler class is $x_{i+1}-x_i-\epsilon$.
Thus in either case
\[
 (\mathcal S_i-\id)\Phi_K(f)
 =\Phi_K\bigl((x_{i+1}-x_i-\epsilon)\partial_i f\bigr)
 =\Phi_K\bigl((\BS_i-\id)f\bigr),
\]
by \eqref{eq:BS-expansion}. Stability of $I_{\rfunc,K}$ follows from $I_{\rfunc,K}=\ker\Phi_K$.
\end{proof}

\begin{corollary}\label{cor:deg_hecke_action_xdelta}
The assignments $s_i\mapsto\mathcal S_i$ for $s_i\in S'$, together with multiplication by
$x_j$ and $\epsilon$, define a graded representation of $\Hdeg'$ on $H^*_K(\SX_{\rfunc};\Z)$.
These representations extend compatibly to any commutative coefficient ring.
In particular, when $\Psi=\Delta^+$, they define an action of the full degenerate affine Hecke algebra $\Hdeg$.
\end{corollary}

\begin{proof}
The relations of the polynomial representation descend through $\Phi_K$
by Theorem~\ref{thm:geometric_hecke}. Coefficient extension follows from
Corollary~\ref{cor:coho_xdelta_R}.
\end{proof}

\subsection{Correspondences realizing the operators}

Let $d=\dim_\C X$, and define the closed subvariety
\[
  Z_{\rfunc}=\{(L_\bullet,L'_\bullet)\in X\times X\mid L_n=L'_n\},
\]
endowed with its reduced subscheme structure.
Because the forgetful morphism $\pi_i$ preserves the terminal lattice $L_n$ (as $i<n$),
all correspondences constructed below are supported on $Z_{\rfunc}$.

When $\epsilon_i-\epsilon_{i+1}\in\Psi$, let $C_i=Y_i\times_{B_i}Y_i$, regarded as a closed subvariety
of $X\times X$ via $j_i\times j_i$.
Since $\dim_\C Y_i=d-1$ and $\dim_\C B_i=d-2$, the fiber product $C_i$ has dimension $d$,
defining an equivariant fundamental class $[C_i]\in H_K^{2d}(X\times X;\Z)$.
We define the correspondence class
\[
  \Gamma_i=[\Delta_X]+[C_i]\in H_K^{2d}(X\times X;\Z).
\]
Let $q_1,q_2:C_i\to Y_i$ denote the two projections.
The fiber square
\[
\begin{tikzcd}
  C_i\arrow[r,"q_2"]\arrow[d,"q_1"'] & Y_i\arrow[d,"\pi_i"] \\
  Y_i\arrow[r,"\pi_i"'] & B_i
\end{tikzcd}
\]
is Cartesian with $\pi_i$ smooth and proper.
By flat base change for fiber integration, we have $q_{1*}q_2^*=\pi_i^*\pi_{i*}$.
Under the correspondence action \eqref{eq:correspondence_action}, the class $[C_i]$ therefore acts by
\[
  j_{i*}q_{1*}q_2^*j_i^*=j_{i*}\pi_i^*\pi_{i*}j_i^*,
\]
so $[\Delta_X]+[C_i]$ acts by \eqref{eq:hecke_case_A}.

When $\epsilon_i-\epsilon_{i+1}\notin\Psi$, the forgetful morphism $\pi_i:X\to B_i$ is defined on the entire variety $X$, and the corresponding
fiber product $C_i=X\times_{B_i}X$ has dimension $d+1$, so its fundamental class $[C_i]\in H_K^{2d-2}(X\times X;\Z)$ has
cohomological degree $2d-2$.
To obtain a degree-zero correspondence, we take the cup product with $\operatorname{pr}_1^*e_K(\mathcal E_i)$:
\[
  \Gamma_i=[\Delta_X]+\operatorname{pr}_1^*e_K(\mathcal E_i)\cup[C_i]\in H_K^{2d}(X\times X;\Z).
\]
By flat base change for $C_i=X\times_{B_i}X$ and the projection formula for $\operatorname{pr}_1$, the action of $\Gamma_i$ coincides with
\eqref{eq:hecke_case_B}.
